% !TEX root = chapter-standalone.tex




\section{Sums}
\label{sec:arithmetic-sums}


\subsection{Commutativity}


The equation $ x + y = y + x$  describes the \emph{commutativity property} of addition.


Let us look at the analogous statement on the level of sets:
\begin{equation}
\setA \setdisunion \setB \isomorphic \setB \setdisunion \setA.
\end{equation}
To prove this statement, we will specify a function $\setA \setdisunion \setB \to \setB \setdisunion \setA$ that is an isomorphism. For this, recall that 
\begin{equation}
\setA \setdisunion \setB = ...
\end{equation}
and
\begin{equation}
\setB \setdisunion \setA = ...
\end{equation}
The only difference between these two sets is that, for $\setA \setdisunion \setB$, the set $\setA$ is marked by the index ``$1$'' as being the first summand and $\setB$ is marked by ``$2$'' as being the second summand, while for $\setB \setdisunion \setA$, the set $\setB$ is marked as first, and $\setA$ is marked as second. We thus define the following function which simply swaps the indices that mark which set is first and which is second: 
\begin{equation}\label{eq:braiding-sum-sets}
\begin{aligned}
        \braiding_{\setA, \setB} \colon & \setA \setdisunion \setB & \to & \ \setB \setdisunion \setA, & \ \\
        & \tup{1, \ela} & \mapsto &\  \tup{2, \ela} & \ \text{ for } \ \tup{1, \ela} \setin \makeset{1} \cartprod \setA
         \\
        & \tup{2, \elb} & \mapsto &\  \tup{1, \elb} & \ \ \text{ for } \ \tup{2, \elb} \setin \makeset{2} \cartprod \setB \\
\end{aligned}
\end{equation}
We call this function $\braiding_{\setA, \setB}$: the letters $\braiding$ stand for ``braiding'' and the indices keep track of which sets are involved and in which order.

Now observe that there is a function which is an obvious candidate to be an inverse of $\braiding_{\setA, \setB}$, namely the same kind of swap again (but starting from $\setB \setdisunion \setA$): 
\begin{equation}\label{eq:braiding-sum-sets-inverse}
\begin{aligned}
        \braiding_{\setB, \setA} \colon & \setB \setdisunion \setA & \to & \ \setA \setdisunion \setB, & \ \\
        & \tup{1, \elb} & \mapsto &\  \tup{2, \elb} & \ \text{ for } \ \tup{1, \elb} \setin \makeset{1} \cartprod \setB
         \\
        & \tup{2, \ela} & \mapsto &\  \tup{1, \ela} & \ \ \text{ for } \ \tup{2, \ela} \setin \makeset{2} \cartprod \setA \\
\end{aligned}
\end{equation}
\begin{exercise}
Check that $\braiding_{\setB, \setA}$ is indeed an inverse to $\braiding_{\setA, \setB}$.
\end{exercise}

\begin{solution}
    According to \cref{def:function-isomorphism}, we have to check the two equations
    \begin{equation}\label{eq:braiding-sum-sets-inverse-equations}
        \braiding_{\setA, \setB} \mthen \braiding_{\setB, \setA} = \catidat{\setA \setdisunion \setB}
        \qqand
        \braiding_{\setB, \setA} \mthen \braiding_{\setA, \setB} = \catidat{\setB \setdisunion \setA}.
    \end{equation}
    Both sides of each of these equations are functions with the same domain and the same codomain, so it suffices to check that they take the same value on every element of their domain.

    For the first equation, every element of $\setA \setdisunion \setB$ is either of the form $\tup{1, \ela} \setin \makeset{1} \cartprod \setA$ or of the form $\tup{2, \elb} \setin \makeset{2} \cartprod \setB$. Using first the definition \cref{eq:braiding-sum-sets} of $\braiding_{\setA, \setB}$ and then the definition \cref{eq:braiding-sum-sets-inverse} of $\braiding_{\setB, \setA}$, we compute
    \begin{equation}
        \begin{aligned}
            \tup{1, \ela} & \ \mapsto \ \tup{2, \ela} & \ \mapsto \ & \tup{1, \ela}, \\
            \tup{2, \elb} & \ \mapsto \ \tup{1, \elb} & \ \mapsto \ & \tup{2, \elb}.
        \end{aligned}
    \end{equation}
    In both cases the composite $\braiding_{\setA, \setB} \mthen \braiding_{\setB, \setA}$ returns the element it started with, which is precisely what the identity function $\catidat{\setA \setdisunion \setB}$ does.

    The second equation is checked in exactly the same way, with the roles of $\setA$ and $\setB$ interchanged: every element of $\setB \setdisunion \setA$ is of the form $\tup{1, \elb} \setin \makeset{1} \cartprod \setB$ or of the form $\tup{2, \ela} \setin \makeset{2} \cartprod \setA$, and
    \begin{equation}
        \begin{aligned}
            \tup{1, \elb} & \ \mapsto \ \tup{2, \elb} & \ \mapsto \ & \tup{1, \elb}, \\
            \tup{2, \ela} & \ \mapsto \ \tup{1, \ela} & \ \mapsto \ & \tup{2, \ela},
        \end{aligned}
    \end{equation}
    so that $\braiding_{\setB, \setA} \mthen \braiding_{\setA, \setB} = \catidat{\setB \setdisunion \setA}$.

    Intuitively: swapping the two index labels twice in a row leaves them unchanged. In particular $\braiding_{\setA, \setB}$ is an isomorphism, which proves that $\setA \setdisunion \setB \isomorphic \setB \setdisunion \setA$.
\end{solution}

\todotext{JL: explain `canonical' here}


\subsection{Associativity}

\todotext{JL: Rewrite this section}

The equation $x + (y + z) = (x + y) + z$ describes the \emph{associativity} property of multiplication and addition of natural numbers.

Correspondingly, there is a canonical isomorphism

\begin{equation}\label{eq:sets-sum-associator}
    \associator \colon (\setA \setdisunion \setB) \setdisunion \setC \to \setA \setdisunion (\setB \setdisunion \setC)
\end{equation}

\begin{exercise}\label{ex:associator-cart-prod-and-disjoint-union}
    Guess how $\associator$ is defined for the sum (disjoint union) of sets, and check that your guess is indeed an isomorphism.
\end{exercise}

\begin{solution}
    \begin{equation}\label{eq:sets-sum-associator-def}
        \begin{aligned}
            \associator \colon & (\setA \setdisunion \setB) \setdisunion \setC  \ \  \to         \ \setA \setdisunion (\setB \setdisunion \setC), \\
                               & \begin{cases}
                                     \tup{1, \tup{1, \ela}} & \ \ \mapsto   \ \tup{1, \ela}         \\
                                     \tup{1, \tup{2, \elb}} & \ \ \mapsto  \ \tup{2, \tup{1, \elb}} \\
                                     \tup{2, \elc}          & \ \ \mapsto \ \tup{2, \tup{2, \elc}}.
                                 \end{cases}
        \end{aligned}
    \end{equation}
\end{solution}


\subsection{The empty set is like the number zero}

Analogous to the equations $0 + x = x$ and $x + 0 = x$, we have
\begin{equation}\label{eq:singleton-set-is-neutral-for-sum}
    \Emptyset \setdisunion \setA \isomorphic \setA \quad \text{and} \quad \setA \setdisunion \Emptyset \isomorphic \setA.
\end{equation}

And here again we have canonical isomorphisms
\begin{equation}\label{eq:leftunitor-sets-sum}
    \defmapcomma{
        \leftunitor
    }{
        \Emptyset \setdisunion \setA
    }{
        \to
    }{
        \setA
    }{
        \tup{2, \ela}
    }{
        \ela
    }
\end{equation}
and
\begin{equation}\label{eq:rightunitor-sets-sum}
    \defmapperiod{
        \rightunitor
    }{
        \setA \setdisunion \Emptyset
    }{
        \to
    }{
        \setA
    }{
        \tup{1, \ela}
    }{
        \ela
    }
\end{equation}
which we also call \emph{left unitor} and \emph{right unitor}, respectively.



\subsection{Arbitrary dependent sums}

There is also an unordered and infinite version of the sum of sets.
Given a family~$\makeset{ \setA_{\setIel} }_{\setIel \setin\setI}$, with arbitrary index set $\setI$, we define the unordered sum as
\begin{equation}
    \underset{\setIel \setin\setI}{\stylesets{\sum}} \setA_{\setIel} \definedas \bigsetunion_{\setIel \setin\setI} \makeset{ \setIel } \cartprod \setA_{\setIel}.
\end{equation}
Also here there is no restriction on the size of $\setI$.
Note that for the unordered version of the sum of sets we have not needed to switch to using functions -- the union of sets is already an ``unordered'' operation.
In the definition of $\setAn{1} \setdisunion \dots \setdisunion \setAn{n}$, the ordering came solely from the ordinals in the index set $\makeset{1, \dots, n}$.


